\documentclass{article}
\usepackage{amsmath,amsthm,amssymb} % Required for math stuff
\usepackage{graphicx} % Required for inserting images
\usepackage{tikz} % Required for drawing graphs
\usepackage{hyperref} % Required for hyperlinks

\newenvironment{conjecture}[2][Conjecture]{\begin{trivlist}
\item[\hskip \labelsep {\bfseries #1}\hskip \labelsep {\bfseries #2.}]}
{\end{trivlist}}

\title{Project 3}
\author{Mengxiang Jiang}
\date{\today}

\begin{document}

\maketitle

\section{Introduction}
One of my favorite websites to visit is the Online Encyclopedia of 
Integer Sequences (OEIS) \cite{sloane2003line}. It is a vast database 
of integer sequences that people from all over the world can contribute 
to and use for research. Each sequence has a unique identifier, 
starting with the letter `A' followed by six digits. For example, 
the sequence of prime numbers is identified as A000040. Today, I want to 
explore the very first sequence in the database: A000001. 
This sequence counts the number of groups of order $n$. 
I will not go into much detail about what a group is, other 
than it is a set of elements with an operation that satisfies 
certain properties. The order of a group is just the number 
of elements in the group. Let us take a look at the first few 
terms of the sequence A000001:
\begin{center}
\begin{tabular}{|c|c|c|c|c|c|c|c|c|c|c|c|}
\hline
$n$ & 0 & 1 & 2 & 3 & 4 & 5 & 6 & 7 & 8 & 9 & 10 \\
\hline
$a(n)$ & 0 & 1 & 1 & 1 & 2 & 1 & 2 & 1 & 5 & 2 & 2 \\
\hline
\end{tabular}
\end{center}
The sequence doesn't look like it's growing very quickly 
at first glance. However, as we look at larger values of $n$, 
we can see that the number of groups of order $n$ does 
indeed increase. For instance, there are 267 groups of order 64. 
So how fast does the sequence grow? This is a surprisingly 
difficult question, and to give you a sense of the complexity, 
how about we play a little game by Matthew Macauley \cite{macauley_2025}. 
Here are the values of the number of groups of order $n$ for $n$ 
from 1014 to 1023:
\begin{center}
\begin{tabular}{|c|c|c|c|c|c|c|c|c|c|c|c|c|}
\hline
$n$ & 1013 & 1014 & 1015 & 1016 & 1017 & 1018 & 1019 
& 1020 & 1021 & 1022 & 1023 \\
\hline
$a(n)$ & 1 & 23 & 2 & 12 & 2 & 2 & 1 & 37 & 1 & 4 & 2\\
\hline
\end{tabular}
\end{center}
Can you guess the number of groups of order 1024? 
Take a moment to think about it before scrolling down for the answer.

\pagebreak

\section{Groups of Order $2^n$ Dominates Conjecture}
So the number of groups of order 1024 is 49,487,367,289 
which was only determined in 2021 by David Burrell \cite{burrell2022number}. 
That's a huge jump from the previous values! One of the reasons why
there are so many groups of order 1024 is that 1024 is a power of 2.
Without going into too much detail, groups of order $p^n$ 
for a prime number $p$ tend to have more complex structures and 
thus more distinct groups.
For example, $n = 729 = 3^6$ has 504 groups. However, powers of 2 are so
dominant that they account for more than 99\% of all groups of order 
less than 2000. Calculating the exact number of groups of 
order $n$ as $n$ grows is still very difficult, 
with John Conway saying that the ``human race will never know the 
exact number of groups of order 2048."
Which brings us to the main conjecture I want to discuss today:
\begin{conjecture}{Groups of Order $2^n$ Dominates}
    Almost all finite groups are of order $2^n$. That is
    \[
        \lim_{n \to \infty} \frac{a(2^n)}{\sum_{k=1}^{n} a(k)} = 1.
    \]
\end{conjecture}
Given the information we have so far, this conjecture seems almost
trivially true. However, it has not been proven rigorously yet.

\begin{section}{Bibliography}
    \bibliographystyle{plain}
    \bibliography{refs}
\end{section}

\end{document}
